\documentclass[10pt,a4paper]{amsart}
\usepackage[margin=19mm]{geometry}
\usepackage{amsmath,amssymb,mathtools}
\usepackage[scr=rsfs,cal=euler]{mathalfa}
\usepackage{xfrac}
\usepackage[unicode,hidelinks,bookmarks=false]{hyperref}
\newcommand{\ie}{i.e.\ }
\newcommand{\cf}{cf.\ }
\newcommand{\resp}{respectively\ }
\newcommand{\Subsection}{\subsection}
\providecommand{\nobreakdash}{\nobreak\hbox{-}}
\setlength{\emergencystretch}{2em}
\multlinegap0pt
\usepackage{bm}
\usepackage{bbm}
\usepackage{dsfont}
\usepackage{xspace}
\usepackage{cancel}
\usepackage{mathtools}
\usepackage{amsthm}
\makeatletter
\@ifundefined{assumption}{\newtheorem{assumption}{Assumption}}{}
\@ifundefined{lemma}{\newtheorem{lemma}{Lemma}}{}
\@ifundefined{theorem}{\newtheorem{theorem}{Theorem}}{}
\makeatother
\title[Accelerating Discrete Diffusion Models with Parallel-In-Time]{Accelerating Discrete Diffusion Models with Parallel-In-Time Sampling}
\author{Yu Yao, Huanjian Zhou, Andi Han, Wei Huang, Masashi Sugiyama}
\date{Source: arXiv:2607.00773v1 (CC BY 4.0)}
\begin{document}
\maketitle

\noindent\emph{Source and licence.} Accelerating Discrete Diffusion Models with Parallel-In-Time Sampling. Version arXiv:2607.00773v1 (CC BY 4.0), distributed under the Creative Commons Attribution 4.0 International licence, as recorded in the arXiv licence metadata for this exact version and shown on the abstract page https://arxiv.org/abs/2607.00773v1. Full rights notice: licenses/arxiv-2607.00773-CC-BY-4.0.txt.\par\smallskip

\noindent\textit{Editorial scope.} Counted unit: the Theorem whose source label is \texttt{thm:picard-endpoint} in the article (source0.tex lines 2079--2094, source label \texttt{thm:picard-endpoint}), delivered together with the article's own complete original proof of it (source0.tex lines 2096--2191, one \texttt{proof} environment of 96 lines). No mathematics is written, repaired or re-derived: statement and proof are the article's own text line for line. Required context retained uncounted: ass:event-switching (lines 1142--1150); lem:triangular-recursion (lines 2031--2039). Every definition, display and label of those lines is retained unchanged. Omitted from this package and not redistributed: 1--1141, 1151--2030, 2040--2078, 2095--2095, 2192--2812. Nothing inside a retained excerpt is omitted or altered: every retained body is delivered exactly as the article's lines are written, so no omission receipt of any kind accompanies this package. Citations: the retained excerpts carry no citation command at all, so no bibliography entry of the article is transcribed and no reference is redistributed. Reference adaptation: none was needed and none was made; every reference inside the delivered unit resolves inside it, so no unresolved reference or citation is produced. Printed slips kept literal: no printed word, symbol, hypothesis or number of the counted statement or of its proof was altered. Layout and class: the article's own class is \texttt{article} with packages including graphicx, wrapfig, caption, xcolor, neurips\_2026, inputenc, fontenc, hyperref, url, booktabs, amsfonts, nicefrac, microtype, enumitem; the wrapper is the lane's pinned \texttt{amsart} editorial wrapper at 10pt on A4, which restates the theorem-like environments the retained text needs and adds the article's own macro definitions that the retained text needs (none), with extra wrapper packages (bm, bbm, dsfont, xspace, cancel, mathtools). The delivered numbering is the unit's own. Assets: no retained span contains a figure environment, a graphics-inclusion command, a PDF-page inclusion, a table or a TikZ picture; no image file is copied into this package, the package declares no asset dependency and no image file is redistributed with it. Licence: version arXiv:2607.00773v1 of this article is distributed under the Creative Commons Attribution 4.0 International licence, as recorded in the arXiv licence metadata for this exact version and shown on the abstract page https://arxiv.org/abs/2607.00773v1. A full rights notice is delivered at licenses/arxiv-2607.00773-CC-BY-4.0.txt. Characters: every retained excerpt is pure ASCII, so the delivered engine, XeLaTeX with its default Latin Modern text font, reports no missing character.
\begin{assumption}[\bf Normalized event switching]
\label{ass:event-switching}
There exists a constant $L_\star>0$ such that, for every fine cell $q$ and all relevant states $z,z'\in\mathcal X_M$,
\[
\mathbb E_{\omega_q}\left[\left|C_q(z,\omega_q)\triangle C_q(z',\omega_q)\right|\right]
\le
L_\star\rho_q d_H(z,z').
\]
\end{assumption}

\begin{lemma}[Triangular Picard recursion]
\label{lem:triangular-recursion}
Under Assumption~\ref{ass:event-switching}, for every $m=0,\ldots,M$ and $k\ge0$,
\[
e_{k+1}(m)
\le
\sum_{r=0}^{m-1}b_r e_k(r).
\]
\end{lemma}

\begin{theorem}[Picard endpoint convergence]
\label{thm:picard-endpoint}
Let
\[
B=\sum_{r=0}^{M-1}b_r,
\qquad
E_0=\max_{0\le m\le M}e_0(m).
\]
Under Assumption~\ref{ass:event-switching}, the block endpoint iterates converge in expected Hamming distance to a limit $Y^{(\infty)}$, and for every $K\ge0$,
\[
\mathbb E\left[d_H\left(Y^{(K)},Y^{(\infty)}\right)\right]
\le
E_0 e^B\frac{B^K}{K!}.
\]
For block $n$, the same statement holds with $B=B_n$. One may always take $E_0\le d$.
\end{theorem}

\begin{proof}
We first prove a bound on adjacent endpoint residuals. We claim that for every $K\ge0$,
\[
e_K(M)\le E_0\frac{B^K}{K!}.
\]
For $K=0$, this is immediate from the definition of $E_0$. Now let $K\ge1$. Repeatedly applying Lemma~\ref{lem:triangular-recursion} gives
\[
e_K(M)
\le
\sum_{r_1<M}b_{r_1}e_{K-1}(r_1).
\]
Applying the recursion to $e_{K-1}(r_1)$ gives
\[
e_{K-1}(r_1)
\le
\sum_{r_2<r_1}b_{r_2}e_{K-2}(r_2).
\]
Substituting,
\[
e_K(M)
\le
\sum_{r_2<r_1<M}b_{r_1}b_{r_2}e_{K-2}(r_2).
\]
Continuing this expansion for $K$ steps yields
\[
e_K(M)
\le
\sum_{r_K<r_{K-1}<\cdots<r_1<M} b_{r_1}\cdots b_{r_K} e_0(r_K).
\]
Since $e_0(r_K)\le E_0$,
\[
e_K(M)
\le
E_0 S_K,
\qquad
S_K:=\sum_{r_K<\cdots<r_1<M} b_{r_1}\cdots b_{r_K}.
\]
It remains to bound $S_K$. Expanding $B^K$ gives
\[
B^K=\left(\sum_{r=0}^{M-1}b_r\right)^K
=\sum_{(i_1,\ldots,i_K)\in\{0,\ldots,M-1\}^K} b_{i_1}\cdots b_{i_K}.
\]
All terms are nonnegative. Hence $B^K$ is at least the sub-sum over tuples with $K$ distinct indices. For each unordered subset $\{a_1,\ldots,a_K\}$ of $K$ distinct indices, the product $\prod_{a}b_a$ appears exactly $K!$ times in the distinct-index tuple expansion, once for each permutation. The strictly decreasing chains in $S_K$ are in one-to-one correspondence with these unordered subsets. Therefore
\[
B^K\ge K!S_K,
\qquad
S_K\le \frac{B^K}{K!}.
\]
Thus
\[
e_K(M)\le E_0\frac{B^K}{K!}.
\]

Now fix $L>K$. By the triangle inequality for Hamming distance,
\[
d_H(Y^{(L)},Y^{(K)})
\le
\sum_{\ell=K}^{L-1}d_H(Y^{(\ell+1)},Y^{(\ell)}).
\]
Taking expectations and using $Y^{(\ell)}=Z_M^{(\ell)}$ gives
\[
\mathbb E[d_H(Y^{(L)},Y^{(K)})]
\le
\sum_{\ell=K}^{L-1}e_\ell(M)
\le
E_0\sum_{\ell=K}^{L-1}\frac{B^\ell}{\ell!}.
\]
The exponential series converges, so the right-hand side tends to zero as $K,L\to\infty$. Hence the endpoints are Cauchy in expected Hamming distance and have a limit, denoted $Y^{(\infty)}$.

Letting $L\to\infty$ in the previous inequality gives
\[
\mathbb E[d_H(Y^{(\infty)},Y^{(K)})]
\le
E_0\sum_{\ell=K}^{\infty}\frac{B^\ell}{\ell!}.
\]
Finally,
\[
\sum_{\ell=K}^{\infty}\frac{B^\ell}{\ell!}
=
\frac{B^K}{K!}\sum_{j=0}^{\infty}\frac{B^j K!}{(K+j)!}.
\]
Since $K!j!\le (K+j)!$, we have $K!/(K+j)!\le1/j!$. Hence
\[
\sum_{j=0}^{\infty}\frac{B^j K!}{(K+j)!}
\le
\sum_{j=0}^{\infty}\frac{B^j}{j!}
=e^B.
\]
Therefore
\[
\mathbb E[d_H(Y^{(\infty)},Y^{(K)})]
\le
E_0 e^B\frac{B^K}{K!}.
\]
This completes the proof.
\end{proof}

\end{document}
